\section{Proposal: \textsc{COMETA-PolJetClass}, an extended JetClass dataset with polarization labels}
\label{sec:proposal}

Building on the open-data discussion of Day~1 and on the deliverable planning
of Day~2, the participants converged on a concrete joint project: the
construction of an extended, \emph{polarization-labeled} JetClass-like dataset
for semi-leptonic diboson processes, together with a community benchmark of the
tagging architectures presented at this workshop. The project deliberately
follows the blueprint of the Particle Transformer paper
(Qu, Li, Qian, arXiv:2202.03772, ICML~2022), which paired a new large-scale
public dataset (JetClass: 100M jets in 10 classes, generated with
MadGraph+Pythia and a Delphes CMS-like detector, anti-$k_t$ $R=0.8$ jets with
$\pT\in[500,1000]$~GeV and per-particle kinematic, identification, and
trajectory-displacement features) with a new architecture (ParT, with pairwise
interaction features in the attention bias) and a standardized evaluation
protocol. The sections below mirror that structure.

\subsection{The \textsc{COMETA-PolJetClass} dataset}

\paragraph{Processes and classes.} The dataset targets semi-leptonic $ZV$
production, $V=W,Z$ decaying hadronically: $pp\to Z(\to\ell\ell)\,W(\to q\bar
q')$ and $pp\to Z(\to\ell\ell)\,Z(\to q\bar q)$, in both QCD-induced and
VBS/EW ($ZVjj$) phase spaces, with the hadronically decaying boson boosted and
captured by a single large-radius jet. The primary jet classes are
$\{W_L,\,W_T,\,Z_L,\,Z_T,\,q/g\}$, the light quark/gluon class being drawn from
$Z(\to\ell\ell)$+jets to match the signal topology. The leptonic $Z$ leg is
retained in the event record so that \emph{production-level} observables
(scattering angle, system kinematics) are available alongside the jet, unlike
in JetClass, which stores isolated jets only.

\paragraph{Improved MC accuracy.} The distinguishing ambition, motivated by the
theory review and the MC roundtable, is generation accuracy beyond the
LO-template standard, in two directions:
\begin{enumerate}
\item \emph{Perturbative accuracy.} Polarized signals are defined in the
double-pole approximation with documented frame conventions, generated at
NLO~QCD accuracy (SHERPA nLO+PS and POWHEG-BOX-RES NLO+PS), supplemented by LO
multi-jet-merged MadGraph samples; fixed-order DPA predictions (NLO and, where
available, NNLO~QCD) are used to validate or reweight the samples. This
addresses the 10--50\% biases of LO polarized templates highlighted in the
experimental review.
\item \emph{Multiplicity of generators and showers.} Every process/class is
produced in several generator$\times$shower combinations (at minimum
MadGraph+Pythia8, SHERPA native, POWHEG+Pythia8, POWHEG+Herwig7, spanning
dipole and angular-ordered showers, with varied hadronization models), and
each event carries a provenance label. Modeling robustness thus becomes a
first-class, \emph{measurable} property of any tagger trained on the dataset,
directly probing the up-to-$\sim$40\% shower/hadronization vulnerability of
low-level taggers reported at the workshop.
\end{enumerate}

\paragraph{Detector emulation, features, and labels.} For comparability with
JetClass, events are passed through the same Delphes CMS-like card, with
anti-$k_t$ $R=0.8$ jets in $\pT\in[400,1000]$~GeV, and per-particle features
form a strict superset of the JetClass ones (kinematics, particle
identification, trajectory displacement). Each jet additionally carries truth
labels: the boson helicity state in the declared frame, the true decay angle
$\costh$ (and azimuth $\varphi^{*}$), the production/scattering angle
$\Theta$ of the diboson system, the truth subjet four-momenta, and the boson
$\pT$. A target size of $\mathcal{O}(10)$M jets per class per
generator/shower combination keeps the total within an order of magnitude of
JetClass while enabling scaling-law studies.

\subsection{Architectures to be compared}

The benchmark will confront, on identical inputs and splits, the families of
approaches presented at the workshop:
(i) a high-level, IRC-safe baseline: BDTs/MLPs on groomed mass,
$\tau_{21}$, $D_2$, energy correlators, and the $p_\theta$ observable of the
De--Rentala--Shepherd method;
(ii) constituent-based deep networks: ParticleNet, a plain Transformer, and
ParT with its pairwise-interaction attention bias, in the five-class
configuration and in the ATLAS-style multi-task configuration with
$\costh$ and subjet-kinematics regression heads;
(iii) Lund-plane approaches: LundNet as classical reference and the LP2B
quantum tree tensor network, testing whether IRC-safe declustering inputs
retain their robustness advantage at scale;
(iv) quantum-informed hybrids: QUIVER-style injection of quantum Fisher
information tokens into ParT;
(v) decorrelated variants: the 3D~DDT applied to the best classifier scores,
flattening the QCD response in mass, $\pT$, and the polarization-sensitive
observable; and
(vi) the amplitude-assisted optimal observable constructed from matrix
elements, serving as the per-event information upper bound against which all
learned approaches are gauged.

\subsection{Benchmark tasks and metrics}

Following the ParT evaluation protocol, but extended to the polarization
context, four tasks are defined:
\begin{enumerate}
\item \emph{Tagging}: each polarized boson class versus $q/g$, reported as
accuracy, AUC, and background rejection at 50\% and 99\% signal efficiency,
\emph{separately per polarization state}, since the workshop showed that
$q/g$ rejection depends on polarization.
\item \emph{Polarization classification}: $V_L$ versus $V_T$ (and the joint
$W$-vs-$Z$ $\times$ $L$-vs-$T$ problem), with AUC complemented by a
physics-level figure of merit: the expected uncertainty on the longitudinal
fraction $\fL$ from template pseudo-fits including the QCD background, which
is what an LHC analysis ultimately optimizes.
\item \emph{Regression}: per-jet estimation of $\costh$ (and optionally
$\varphi^{*}$ and the subjet kinematics), evaluated through resolution
($\sigma_{68}$), bias versus true $\costh$, and---critically---response
uniformity between $L$ and $T$ samples, so that the regressed observable can
be used in an unbiased template fit.
\item \emph{Decay versus production angles}: a dedicated study quantifying
whether the decay angle $\costh$ (jet-internal information) or the production
angle $\Theta$ (event-level information, accessible thanks to the retained
leptonic leg) carries more polarization sensitivity per event, using
Fisher-information and mutual-information estimates as well as the degradation
of the $\fL$ pseudo-fit when either input is removed. This directly informs
whether future analyses should invest in jet-substructure regression, in
event-level reconstruction, or in their combination.
\end{enumerate}
All results are to be reported with a robustness matrix (train on
generator/shower $A$, test on $B$) and, for decorrelated taggers, with the
sculpting of the QCD templates quantified by a Jensen--Shannon divergence.

\subsection{Suggestions}

\textcolor{red}{%
The following suggestions were added by the author of this summary.
(1)~Store, for every event, the per-helicity DPA matrix-element weights
(including the interference term): this makes the amplitude-assisted optimal
observable computable on the released events, allows reweighting between
polarization definitions and reference frames without regeneration, and turns
the definition dependence itself into a quantifiable systematic.
(2)~Provide helicity labels in at least two frames (e.g.\ the $VV$
center-of-mass helicity frame and the Collins--Soper frame), since polarization
fractions are frame-dependent and ATLAS/CMS conventions differ.
(3)~Include an explicit ``interference/off-shell'' event category rather than
forcing every event into a pure-helicity class, so benchmark fits can propagate
the non-closure of the polarized decomposition.
(4)~Reweight all classes to a common boson-$\pT$ spectrum (or provide the
weights) to prevent taggers from learning production kinematics instead of
decay structure---particularly important for the decay-vs-production-angle
study of Task~4.
(5)~Publish fixed train/validation/test splits, evaluation code, and a public
leaderboard in the style of the \texttt{jet-universe} repository, with a
versioned DOI (Zenodo/CERN Open Data), and require metrics to be quoted as
mean~$\pm$~spread over at least five trainings.
(6)~Add a small full-simulation companion sample (ATLAS/CMS open data) to test
ParT-style fine-tuning and quantify the Delphes-to-full-sim transfer gap before
any in-experiment deployment.
(7)~Extend the ParT pairwise interaction features with polarization-motivated
pair variables (e.g.\ the momentum-sharing $z$ and pair mass in the jet rest
frame), a cheap architectural handle that directly encodes the physics of the
$L$/$T$ decay distributions~$\propto\sin^2\theta^{*}$ vs
$(1\pm\cos\theta^{*})^2/2$.
(8)~Include the azimuthal decay angle $\varphi^{*}$ among the regression
targets from the start: it gives access to the off-diagonal spin-density-matrix
elements and thus connects this dataset to the quantum-information program
(entanglement observables) highlighted in the experimental review.}

\section*{Acknowledgements}
This summary is based on the material uploaded to the workshop Indico page
(\url{https://indico.cern.ch/event/1601346/}). The workshop was supported by
the COMETA COST Action CA22130.
